% MATH Companion Part VI / VI-29 / explanatory bridge only; no exercise replacement.
\paragraph{What is actually counted by Krull dimension.}
A scheme can have far more points than one can draw on a coordinate grid.
Some points correspond to entire irreducible algebraic subsets, and the
specialization relation orders these points.  Krull dimension measures the
length of the longest possible sequence of strict specializations, not the
number of closed points or the number of generators of a ring.

For a commutative ring \(A\), a chain
\[
 \mathfrak p_0\subsetneq\mathfrak p_1\subsetneq\cdots
 \subsetneq\mathfrak p_r
\]
of prime ideals has \emph{length} \(r\): there are \(r\) strict
inclusions, although there are \(r+1\) prime ideals.  The Krull
dimension \(\dim A\) is the supremum of these lengths.  This
convention removes a common off-by-one error in computations.

\paragraph{Why the topology reverses the chain.}
The closed subset corresponding to \(\mathfrak p\) is
\(V(\mathfrak p)\), and it is irreducible because \(\mathfrak p\)
is prime.  Inclusion of ideals reverses inclusion of vanishing sets:
\[
 \mathfrak p\subseteq\mathfrak q
 \quad\Longleftrightarrow\quad
 V(\mathfrak q)\subseteq V(\mathfrak p).
\]
The reverse implication is valid \emph{for prime ideals} because
\(I(V(\mathfrak p))=\sqrt{\mathfrak p}=\mathfrak p\).
Thus the algebraic chain becomes a chain of increasingly smaller
irreducible closed subsets.  Dimension may consequently be computed
in the algebra or in the topology of \(\operatorname{Spec}A\).

\paragraph{Worked calculation: an affine plane.}
Let \(A=k[x,y]\), with \(k\) a field.  Consider
\[
 (0)\subsetneq(x)\subsetneq(x,y).
\]
All three ideals are prime: their quotients are, respectively,
\(k[x,y]\), \(k[y]\), and \(k\), each an integral domain.
The chain therefore proves \(\dim A\ge 2\).  The polynomial-ring
dimension theorem gives \(\dim k[x,y]\le 2\), so
\[
 \boxed{\dim k[x,y]=2.}
\]
Geometrically the reversed chain reads
\[
 \{(0,0)\}=V(x,y)\subsetneq V(x)\subsetneq
 V(0)=\mathbb A^2_k.
\]
It consists of a point on a line in the entire affine plane.
The closed-point coordinates are only a convenient picture; the
generic points of the plane and the line are indispensable in the
scheme-theoretic chain.

\paragraph{Two comparisons that clarify the definition.}
For a field \(k\), the only prime ideal is \((0)\), and therefore
\(\dim k=0\).  For \(k[t]\), the chain
\((0)\subsetneq(t)\) has length one.  Every nonzero prime of
\(k[t]\) is maximal because \(k[t]\) is a principal ideal domain,
so there is no chain of length two and \(\dim k[t]=1\).

Nilpotents do not increase Krull dimension.  For example, set
\(R=k[t,\varepsilon]/(\varepsilon^2)\).  Every prime of \(R\)
contains the nilpotent \(\varepsilon\), and passage to the reduced
quotient gives a bijection preserving prime inclusions:
\[
 \operatorname{Spec}R\simeq
 \operatorname{Spec}(R/(\varepsilon))=
 \operatorname{Spec}k[t],
 \qquad \dim R=1.
\]
The topological space, and hence its Krull dimension, does not see
the nilpotent thickening.  The structure sheaf does see it.

\paragraph{Localization measures dimension near one point.}
Fix \(\mathfrak p\in\operatorname{Spec}A\).  Primes of the local
ring \(A_{\mathfrak p}\) correspond to prime ideals of \(A\)
contained in \(\mathfrak p\).  Therefore
\[
 \dim A_{\mathfrak p}=\operatorname{ht}(\mathfrak p),
\]
where the height counts chains of primes \emph{ending at}
\(\mathfrak p\).  For the plane, the prime \((x)\) has height
one, whereas the closed point \((x,y)\) has height two.  This is
the algebraic foundation for the local-dimension and codimension
calculations in the next sections.
